\section{Introduction}

A tree index answers a query by descending from the root, and the cost of that
descent has two parts: the number of nodes visited, and the work done inside
each. Both are governed by fan-out. Fan-out sets the height of the tree, so it
sets how many nodes a descent visits; and it sets how many entries a visited node
contains, so it sets the work done there.

Fan-out, in turn, is arithmetic on bytes:
%
\begin{equation}
\label{eq:fanout}
f = \left\lfloor \frac{B - h}{k + p + \sigma} \right\rfloor ,
\end{equation}
%
for page size $B$, page header $h$, key size $k$, downlink size $p$ and
per-entry overhead $\sigma$. Every byte an entry does not need is paid for at
every level of every descent, for the life of the index.

In an \emph{extensible} framework the accounting is worse than in a
purpose-built index, because the core does not know what a key means and must
therefore be conservative about it. This paper is about three places where that
conservatism cost bytes or cycles for nothing, in a generalized search
tree~\cite{hellerstein1995generalized} as implemented in PostgreSQL.

The first is not an algorithm but an interface. The framework obliged every
operator class to supply two functions converting between the indexed value and
its stored representation, even when the two coincide and both functions were
identities. Making them optional removes an identity call performed once per
entry of every visited node. The more interesting consequence is that the
absence of a function becomes information: if no conversion is supplied, the core
knows the stored form \emph{is} the original value, and can answer a query from
the index alone~--- something that previously required supplying yet another
function.

The second separates the attributes a query needs to \emph{answer} from those it
needs to \emph{navigate}. Storing the former only in leaves keeps them out of the
keys that set fan-out at every internal level.

The third concerns entries that are modified rather than added. A descent widens
the keys on its path; the original implementation performed each widening as a
deletion followed by an append, moving the data of the page and destroying the
correspondence between entry order and physical order.

\paragraph{Contributions.}

\begin{enumerate}
\item Optional support functions for key conversion, and the reading of their
      absence as a statement about the stored representation, which yields
      index-only scans without an additional function
      (Section~\ref{sec:optional}).
\item Covering attributes for a framework in which the covered column may have
      no operator class of its own, with the fan-out argument that distinguishes
      this case from the corresponding B-tree feature (Section~\ref{sec:include}).
\item In-place overwriting of an entry, which for the common case of a key
      widened during descent moves no data (Section~\ref{sec:overwrite}), and an
      account of the compatibility consequence it surfaced
      (Section~\ref{sec:bytes}).
\end{enumerate}

\section{Background}
\label{sec:background}

A generalized search tree is a balanced tree whose node occupies one page and
whose entries pair a key with a downlink, the key covering everything in the
subtree below. The key type and its operations belong to an operator class
supplied outside the core: \code{consistent} decides whether a subtree may hold
an answer, \code{union} produces a key covering a set of keys, \code{penalty}
prices the widening of a key, \code{picksplit} partitions an overflowing node.
Substituting bounding rectangles yields an R-tree, signatures a full-text index,
ranges a range index.

Two further functions concern representation rather than search.
\code{compress} converted an indexed value into the form stored in the tree,
\code{decompress} converted it back. The intent was to store something smaller
than the value~--- a bounding rectangle in place of a polygon. A third,
\code{fetch}, reconstructs the original value from the stored form and exists so
that a query can be answered without visiting the table.

Height follows from fan-out as $H \approx \log_f N$, so~\eqref{eq:fanout} enters
the cost of a point query as
$\sum_l \omega_l (c_r + f c_c)$~--- a page access per visited node plus a
per-entry cost inside it. Fan-out therefore acts in two directions at once: it
reduces the number of nodes visited and increases the work in each. The changes
below all move $k$ or $\sigma$, that is the denominator of~\eqref{eq:fanout}.

\section{The absence of a function as information}
\label{sec:optional}

\subsection{The problem}

For a substantial share of operator classes no conversion is wanted: the stored
representation is the indexed value. The framework nevertheless required both
functions, so such classes defined identities. Beyond the redundant code this had
two effects.

The core could not tell whether the stored form equalled the original value.
Lacking that knowledge it could not answer a query from the index alone, and an
operator class wanting index-only scans had to supply \code{fetch}~--- a third
function, whose body for these classes is again an identity.

And the identity \code{decompress} was called for every entry of every scanned
page, that is $O(f)$ times per visited node. In terms of
Section~\ref{sec:background} this is a contribution to $c_c$, the per-entry cost,
which by~\eqref{eq:fanout} grows as fan-out improves. An identity function call
is cheap; multiplied by fan-out and by tree height and by every query, it is not
free, and it is pure overhead.

\subsection{The change}

Both functions were made optional, and their absence given meaning. When no
\code{compress} is defined, the core knows the stored representation is the
indexed value itself, and therefore

\begin{itemize}
\item answers queries from the index alone without requiring \code{fetch}; and
\item performs no conversion when scanning a node.
\end{itemize}

Implemented in \code{d3a4f89d8a3}, released in PostgreSQL~11.

\subsection{Why this is worth reporting as a result}

The performance part is a small constant. The part we consider more useful is the
interface part, and it generalizes past this framework.

An extensible system defines what an extension must provide. The obligation is
usually read as a floor: here is the minimum you must implement. But every
mandatory element also destroys information, because a requirement satisfied by
everyone distinguishes nobody. When identity conversions are mandatory, the class
that needs no conversion is indistinguishable from the class that needs one, and
the core must assume the harder case for all. Making the element optional
restores the distinction and lets the core specialize.

This is why the change removes an obligation and \emph{adds} a capability at the
same time, which is otherwise a surprising combination. The cost of a mandatory
interface element is not only the effort of implementing it; it is the loss of
whatever could have been inferred from choosing not to.

There is a practical corollary for extensibility, which is the property this
framework exists to provide. The barrier to writing a new operator class is set
by the volume of obligatory scaffolding, and two of the required functions here
were scaffolding for the majority of classes.

\section{Covering attributes}
\label{sec:include}

\subsection{The problem}

A query can be answered from the index alone when every column it needs is in the
index. The naive route is to add the columns to the key. For a generalized search
tree this fails twice.

First, a key attribute needs an operator class for its type, and the column one
wants in the answer may have none~--- there is no meaningful notion of a
bounding key for it, and no reason there should be. Second, and quantitatively,
a key attribute widens $k$ at \emph{every} level. By~\eqref{eq:fanout} this
lowers fan-out at every internal level, and hence raises $H$; a column that
navigation never consults would make every descent longer.

\subsection{The change}

Index attributes were divided into key and covering. Covering attributes are
stored only in leaf pages, take no part in navigation, do not enter the keys of
internal nodes, and require no operator class for their type. Fan-out at internal
levels is then given by~\eqref{eq:fanout} with $k$ over the key attributes alone,
and at the leaf level with $k + k_{\mathrm{inc}}$. Since height depends on
internal fan-out, covering attributes do not increase it.

Implemented in \code{f2e403803fe}, released in PostgreSQL~12; the corresponding
change for the space-partitioned tree is \code{09c1c6ab4bc}, released in~14.

\subsection{How this differs from the B-tree case}

The feature has a B-tree namesake and a different motivation. There, covering
attributes serve mainly to place a uniqueness constraint on part of an index's
attributes while keeping other columns available in the index; the type of a
covered column is not the difficulty, since a B-tree operator class exists for
essentially every scalar type one would index.

Here the primary case is the column whose type has \emph{no} generalized search
tree operator class, and could not sensibly have one. A spatial index on a
geometry column that must also return a name or an identifier is not an exotic
requirement, and before this change it forced a visit to the table for every
matching row. The two features share a syntax and solve different problems, which
is worth stating because the shared name invites the assumption that the
reasoning transfers.

\section{Modifying an entry in place}
\label{sec:overwrite}

Insertion widens keys along its path: the parent key must cover the newly
inserted entry. Each widening modifies an entry already on a page.

The original implementation performed the modification as a deletion followed by
an append. Entries are stored compactly, so deletion moves all data below the
removed entry, and the modified entry lands at the end of the page. The second
effect is the less obvious one: entry order and physical order stop
corresponding, which costs processor cache locality on every subsequent scan of
that page~--- again a contribution to $c_c$.

The change introduces overwriting in place. When the new entry has the size of
the old one, which is the typical case for a widened key, no data moves at all;
when the sizes differ, only the difference moves. Entry order is preserved.

Implemented in \code{b1328d78f88}, released in PostgreSQL~10. This result was
published previously~\cite{borodin2017memory}, with measurements against
release~10 showing insertion of sorted data about three times faster and about
15\,\% on random data, alongside a smaller index; we restate it here for
completeness of the fan-out picture rather than as a contribution of this paper,
and Section~\ref{sec:eval} re-measures it on a current release.

\section{A page-layout change is a compatibility change}
\label{sec:bytes}

One consequence surfaced during review and deserves its own section, because it
is not specific to this optimisation.

After the change, replaying the write-ahead log produces a page whose bytes
differ from those produced before, although the page is logically equivalent:
the same entries, in the same order, at different offsets. Nothing in the
database's own operation depends on the difference. But tools that verify replay
by comparing a replica's page against the primary's byte for byte do depend on
it, and such tools are how one gains confidence that replay is correct.

The general form: an optimisation of page layout is a change to an interface that
is not written down anywhere. The set of byte-level representations a page may
legitimately have is a contract with everything outside the database that reads
pages, and it has no declaration to consult. We note it here as a hazard for
anyone proposing a similar change, and as an argument for the structural
verification approach we develop elsewhere: comparing pages byte for byte is a
brittle proxy for the property actually wanted, which is that the two pages carry
the same entries.

\section{Evaluation}
\label{sec:eval}

\subsection{Covering attributes}

One million points with an additional integer column, 2000 range queries,
PostgreSQL~18.4.

\begin{table}[t]
\centering
\caption{Three ways to make a column available to a query}
\label{tab:include}
\begin{tabular}{lrrrr}
\toprule
column is & index pages & internal fan-out & accesses/query \\
\midrule
absent & 8\,160 & 122.6 & 4.19 \\
covering & 9\,622 & 123.9 & 4.46 \\
in the key & \multicolumn{3}{c}{\emph{index cannot be created}} \\
\bottomrule
\end{tabular}
\end{table}

The third row is the more direct confirmation of the argument in
Section~\ref{sec:include}, and it is categorical rather than quantitative:
\code{CREATE INDEX \dots{} USING gist(p, i)} fails, because no generalized search
tree operator class exists for the integer type. Putting the column in the key is
not an inferior option here; it is not an option, unless an extension supplying
such a class is installed first. This is the case we claimed distinguishes this
feature from its B-tree namesake, and it is the ordinary case for a spatial index
that must also return an identifier.

The quantitative claim holds as well: internal fan-out is unchanged, 122.6
against 123.9, a difference within rounding. The covering attribute enlarges the
index by 18\,\% and raises accesses per query by 6\,\%, and does so entirely at
the leaf level; the length of a descent is unaffected, which is what
Section~\ref{sec:include} predicted from~\eqref{eq:fanout}.

What is missing from this table is the cost of the alternative when it is
possible. With an extension providing an integer operator class the index in the
key can be built, and the fan-out penalty at internal levels can be measured
rather than merely predicted. We have not done so.

\subsection{Two measurements that did not come out}

We report these because omitting them would misrepresent the state of the
evidence.

\paragraph{Optional support functions.} The claim of
Section~\ref{sec:optional} is a capability, not a speed: an operator class
omitting the conversion functions should be usable, and should support
index-only scans without a fetch function. Our test created such a class on
builds either side of the change and found it accepted by both, which means the
test did not exercise the check it was aimed at rather than that the change did
nothing. The proper test builds an index with such a class and inspects whether
an index-only plan is chosen. The capability claim in
Section~\ref{sec:optional} therefore rests on the code and the documentation, not
on a measurement of ours.

\paragraph{In-place overwrite.} Measured on one million points, three runs, the
change is not distinguishable on sorted insertion (9.72--9.76~seconds on both
sides) and worth about 1.3\,\% on random insertion, with the index marginally
larger rather than smaller. The previously published
figures~\cite{borodin2017memory} for this change were about threefold on sorted
data and about 15\,\% on random, with a smaller index.

We do not conclude that those figures were wrong, because a simpler explanation
is available and we have not excluded it. A \code{point} key occupies 16~bytes,
so overwriting one moves almost nothing, whereas the benefit is proportional to
the volume of data shifted; and at 9.7~microseconds per row the measurement is
dominated by write-ahead logging and by the table, not by index page layout. The
measurement that would settle it uses a large key type and separates index
insertion from logging. Until it is done, the claim in
Section~\ref{sec:overwrite} should be read as it is stated there~--- as prior
published work, on a release of that era, not re-established here.

\section{Limitations}

The three changes act on different terms and are not interchangeable: the first
and third reduce $c_c$, the second protects $k$ at internal levels, and none of
them affects the overlap of keys, which for a multidimensional index dominates
the cost of a selective query. Fan-out is a lever with a limit: raising it
reduces the number of nodes visited but increases the work inside each, so a
larger page is not uniformly better, and past some point only an in-page
structure that avoids examining all $f$ entries can help further.

The interface argument of Section~\ref{sec:optional} rests on one framework and
one pair of functions. We believe it generalizes~--- a mandatory element that
everyone satisfies identically carries no information~--- but we have one
instance of it, not a survey.

\section{Conclusion}

Fan-out is usually treated as a property of the data: keys are the size they are.
In an extensible framework a noticeable part of it is instead a property of the
interface between the core and the code that supplies key semantics, and can be
recovered by changing that interface rather than the data.

Of the three changes reported here, the one we would highlight is the smallest in
its performance effect. Making two support functions optional removes an
obligation from every operator class author and simultaneously grants the core a
capability it did not have, because a requirement that everyone satisfies
identically conveys nothing, and dropping it lets the core learn something from
the choice not to implement. That reading~--- optionality as a channel of
information rather than as a convenience~--- seems to us the transferable part of
this work.
